\section{Preliminaries}
\label{sec:preliminaries}

This section defines notation and gives interface-level definitions for the lattice primitives, the NIZK model, the
blind-signature syntax, ARC syntax and security goals, and pseudorandom functions.



\subsection{Notation}
Let $\lambda$ denote the security parameter and  \(\negl(\lambda)\) be any positive function negligible in  \(\lambda\). We write PPT for
probabilistic algorithms whose runtime is polynomial in \(\lambda\). For integers \(a<b\) we write
\([a,b]\coloneqq\{a,a+1,\ldots,b\}\) and we use \([b]\) as shorthand for \([1,b]\). Throughout this work, we denote vectors in lower-case bold letters such as $\vectorify{b}$ and matrices in upper-case bold letters such as $\mA$.

We work over the cyclotomic ring \(\R=\Z[X]/(X^N+1)\) for power-of-two \(N\) and its quotient \(\Rq=\R/q\R\), identified
with polynomials of degree \(<N\) with coefficients in \(\Z_q\). For \(a\in\Z\), we write \(\langle a\rangle_q\in\{0,\dots,q-1\}\)
for reduction modulo \(q\). When considering $\Z_q$-elements from an interval \([-a,b]\), we identify these with their representatives in \(\{-q/2,\dots,q/2\}\).
For elements $x\in R_q$ we write $||x||\leq B$ to mean that each coefficient of $x$ comes from $[-B,B]$. For vectors $\vectorify{x}\in R_q^\ell$, $||\vectorify{x}||\leq B$ means that the statement holds for all elements of $\vectorify{x}$.

All our constructions share the same modulus $q$, ring $R_q$, and coefficient bound \(\BoundMsg\). We write
\(\mathcal{M}\) for the bounded message space appropriate to the primitive under discussion. In the blind-signature and
ARC layers, the full signed message has length \(\ell_\mu=\ell_m+\ell_k\) and is written as
\(\vecmu=(\vecm\Vert \veck)\). We assume that all public dimension and modulus parameters are collected in
\(\PublicParamsGen\) and honestly generated by \(\GenSetup(1^\lambda)\). All other setup algorithms for primitives obtain
\(1^\lambda,\PublicParamsGen\) as input, and we do not mention this explicitly.
Whenever we use the CRT decomposition \(\Rq\cong \prod_i F_i\), we write
\[
(x_i)_i \Had (y_i)_i := (x_i y_i)_i,
\qquad
(x_i)_i \Hdiv (y_i)_i := (x_i y_i^{-1})_i
\]
whenever every \(y_i\neq 0\). Thus \(\Had\) and \(\Hdiv\) denote slotwise multiplication and, where defined, slotwise
division in the CRT basis.

\subsection{Commitment schemes and aligned linear commitments}
\label{subsec:prelim-commit}
\label{subsec:commitments}

\begin{definition}[Commitment scheme]\label{def:commitment-scheme}
A (non-interactive) commitment scheme is a tuple of PPT algorithms
\[
  \ComScheme=(\ComSetup,\Commit,\ComVerify)
\]
with the following syntax:
\begin{itemize}[leftmargin=1.25em]
  \item $\ComPublicParam\gets \ComSetup(\SecParam,\PublicParamsGen)$ is a PPT algorithm which outputs public parameters $\ComPublicParam$. It is implicit input to all other algorithms.
  \item $c \gets \Commit(m)$ is a PPT algorithm which takes a message $m\in\mathcal{M}$ as input and outputs a commitment $c$ and decommitment $d$.
  \item $\ComVerify(c,m,d)\in\{0,1\}$ deterministically checks validity of the opening.
\end{itemize}
Correctness holds if for all $\ComPublicParam\gets \ComSetup$,  $m\in \mathcal{M}$ and $(c,d)\gets \Commit(m)$ it holds that
\[
  \ComVerify(c,m,d)=1.
\]

The scheme is \emph{computationally binding} if for all PPT adversaries $\Adv$,
\[
  \Pr\Big[\begin{matrix}
  	\ComPublicParam \gets \ComSetup(\SecParam,\PublicParamsGen) \wedge (c,m,d,m',d')\gets \Adv(\ComPublicParam) \wedge  \\
  	m\neq m' \wedge \ComVerify(c,m,d)=1 \wedge \ComVerify(c,m',d')=1
  \end{matrix}
   \Big]\le \negl(\lambda),
\]
where the probability is taken over the randomness of $\ComSetup,\Adv$.

The scheme is \emph{(computationally) hiding} if for all PPT distinguishers $\mathcal{D}$, all 
messages $m_0,m_1\in\mathcal{M}$ and $\ComPublicParam\gets \ComSetup(\SecParam,\PublicParamsGen)$,
\[
  \Big|\Pr[\mathcal{D}^{\mathcal{O}_{\mathsf{Com}}(m_0)}(\ComPublicParam)=1]
  -\Pr[\mathcal{D}^{\mathcal{O}_{\mathsf{Com}}(m_1)}(\ComPublicParam)=1]\Big|\le \negl(\lambda)
\]
where $\mathcal{O}_{\mathsf{Com}}(m_b)$ runs $(c,d)\gets \Commit(m_b)$ and outputs $c$ and the probability is taken over the randomness of $\Commit,\mathcal{D}$.
\end{definition}

\subsubsection{Aligned linear commitment}
We use a bounded linear commitment over \(\Rq\) whose opening is split into a message part and a linear-side
randomness pair. Fix message length \(\ell_\mu\), linear-secret dimension \(\ell_s\), output length
\(\ell_{\ComScheme}\), coefficient bound \(\BoundMsg\), and bounded-support sampling distributions
\(\chi_s,\chi_e\) supported on coefficient intervals inside \([-\BoundMsg,\BoundMsg]\). The public commitment key is
written as
\[
  \mACom=\bigl[\mA_\mu \mid \mA_s \mid I_{\ell_{\ComScheme}}\bigr],
\]
where \(\mA_\mu\getsunif \Rq^{\ell_{\ComScheme}\times \ell_\mu}\) and
\(\mA_s\getsunif \Rq^{\ell_{\ComScheme}\times \ell_s}\).
\begin{itemize}
  \item \(\ComSetup(\SecParam,\PublicParamsGen)\) samples \(\mA_\mu,\mA_s\) as above and outputs
  \(\ComPublicParam:=(\mA_\mu,\mA_s,\BoundMsg)\).
  \item The message space is
  \[
    \mathcal{M}:=\{\vectorify{\mu}\in\Rq^{\ell_\mu}\mid \|\vectorify{\mu}\|_\infty\le \BoundMsg\}.
  \]
  \item \(\Commit(\vectorify{\mu})\) samples \(S\gets \chi_s^{\ell_s}\) and
  \(E\gets \chi_e^{\ell_{\ComScheme}}\) and outputs
  \[
    \vectorify{c}=\mA_\mu\vectorify{\mu}+\mA_s S + E
  \]
  together with opening \((S,E)\).
  \item \(\ComVerify(\vectorify{c},\vectorify{\mu},(S,E))\) accepts iff
  \[
    \vectorify{c}\in\Rq^{\ell_{\ComScheme}},\qquad
    \vectorify{\mu}\in\mathcal{M},\qquad
    \|S\|_\infty,\|E\|_\infty\le \BoundMsg,
  \]
  and \(\vectorify{c}=\mA_\mu\vectorify{\mu}+\mA_s S + E\).
\end{itemize}

\begin{definition}[Module-SIS and Module-ISIS over \(R_q\)]\label{def:sis-isis}
Fix integers \(n,m\ge 1\), modulus \(q\ge 2\), and bound \(\beta>0\).
\begin{itemize}[leftmargin=1.25em]
  \item \textbf{\(\mathsf{M-SIS}_\infty(n,m,q,\beta)\).} Given \(\mA\in R_q^{n\times m}\), find a non-zero
  \(\vectorify{s}\in R^m\) such that \(\mA\vectorify{s}\equiv 0 \pmod q\) and \(\|\vectorify{s}\|_\infty\le \beta\).
  \item \textbf{\(\mathsf{M-ISIS}_\infty(n,m,q,\beta)\).} Given \((\mA,\vect)\in R_q^{n\times m}\times R_q^n\), find
  \(\vectorify{s}\in R^m\) such that \(\mA\vectorify{s}\equiv \vect \pmod q\) and
  \(\|\vectorify{s}\|_\infty\le \beta\).
\end{itemize}
\end{definition}

\begin{lemma}[Binding of the aligned commitment]\label{lem:ajtai-binding}
Let $\Adv$ be any PPT adversary that breaks the binding property of the aligned commitment scheme with
probability $p$ for random matrices $\mA_\mu,\mA_s$ and coefficient bound $\BoundMsg$. Then there exists a PPT
attacker $\Adv'$ that breaks the $\mathsf{MSIS}_\infty(n,m,q,2\BoundMsg)$ assumption with essentially the same
probability and runtime.
\end{lemma}

\begin{proof}
If two distinct openings \((\vectorify{\mu},S,E)\neq(\vectorify{\mu}',S',E')\) verify the same commitment,
then
\[
\mA_\mu(\vectorify{\mu}-\vectorify{\mu}')+\mA_s(S-S')+(E-E')=0
\]
in \(\Rq^{\ell_{\ComScheme}}\). The concatenated difference is non-zero and coefficient-wise bounded by
\(2\BoundMsg\), hence yields an \(\mathsf{MSIS}_\infty\) solution with essentially the same probability and runtime.
\end{proof}

The hiding property of the aligned commitment is stated below together with the aligned
\(h_{\mathrm{tran}}\) target used by the issuance protocol.

\subsection{Lattice trapdoors and rational hashing}
\label{subsec:prelim-vsis}
We recall only the definitions needed to state our construction, where we follow the definitions \cite{EPRINT:DKLW25}. We defer additional definitions such as the $\mathsf{vSIS}$- and $\mathsf{GenISIS}_f$-assumptions to Appendix \ref{app:assumptions lattices}. 

\begin{definition}[Lattice trapdoor suite]\label{def:trapdoor-suite}
A lattice trapdoor suite over \(\Rq\) consists of PPT algorithms
$  (\Setup,\TrapGen,\SampPre)
$
 such that:
\begin{itemize}[leftmargin=1.25em]
\item $(\PublicParamsTrapdoor)\gets \Setup(\SecParam,\PublicParamsGen)$ generates admissible public parameters \((n,m,q,\BoundSig)\) contained in $\PublicParamsTrapdoor$.
  \item \((\mA,\trapdoor)\gets \TrapGen(\PublicParamsTrapdoor)\) outputs a public matrix \(\mA\in\Rq^{n\times m}\) and a trapdoor
  \(\trapdoor\), where the distribution of \(\mA\) is statistically close to uniform over \(\Rq^{n\times m}\).
  \item On input \((\PublicParamsTrapdoor,\trapdoor,\vect)\) with \(\vect\in\Rq^n\), the algorithm \(\SampPre(\trapdoor,\vect)\) outputs
  \(\vecu\in R^m\) such that \(\mA\vecu=\vect\) in \(\Rq^n\), except with negligible failure probability, and
  \(\|\vecu\|_\infty\le \BoundSig\) except with negligible probability.
\end{itemize}
\end{definition}
We defer details of the sampler that we use in this work to Appendix~\ref{app:gaussian-sampler}.

\begin{definition}[Trapdoor rational hash]\label{def:trapdoor-rational-hash}
A trapdoor rational hash scheme over \(\Rq\) is a tuple
\[
  \TDRH=(\Setup,\TrapGen,\Eval,\SampPre)
\]
consisting of the lattice trapdoor suite above together with a public evaluation algorithm. We extend
\(\Setup\) so that $\PublicParamsTrapdoor$ also contains a public hash description \(B\), a randomness space \(\mathcal{X}\) and an
admissibility error bound \(\delta\). For fixed \(B\), the algorithm
\(\Eval(\PublicParamsTrapdoor,m,\chi)\) takes a message \(m\in\mathcal{M}\) from the ambient message space fixed by
\(\PublicParamsGen\) and witness randomness \(\chi\in\mathcal{X}\) and outputs a
target \(\vect\in\Rq^n\) or \(\bot\) if the input is inadmissible.
 For a public key \(B\) contained in $\PublicParamsTrapdoor$, define the admissible message subspace
\[
  \mathcal{M}_B:=\left\{m\in\mathcal{M} :
  \Pr_{\chi\gets \mathcal{X}}\big[\Eval(\PublicParamsTrapdoor,m,\chi)=\bot\big]\le \delta\right\}\text{.}
\]
Correctness requires that for
every \(m\in\mathcal{M}_B\) and every \(\chi\in\mathcal{X}\) with \(\Eval(\PublicParamsTrapdoor,m,\chi)=\vect\neq\bot\), the output
\(\vecu\gets \SampPre(\PublicParamsTrapdoor,\trapdoor,\vect)\) satisfies \(\mA\vecu=\vect\) and \(\|\vecu\|_\infty\le \BoundSig\), except with
negligible probability.
\end{definition}
This interface does not by itself require the public target \(\vect\) to hide \(m\); in our aligned instantiation,
that privacy is supplied later by the hiding statements of Section~\ref{subsec:aligned-hiding}.

In this work, we use the BB-tran-type rational hash defined in \cite[Table~1]{EPRINT:DKLW25}. In the concrete
ARC-SPRUCE instantiation below, we specialize the generic target space to the single-coordinate case \(n=1\).

\begin{definition}[Rational hash, admissibility, and eliminated form]\label{def:rational-hash}
Our concrete instantiation uses a public description \(B=(B_0,B_1,B_2,B_3)\), where \(B_0,B_3\in\Rq\) and
\(B_1,B_2\) denote public linear maps into \(\Rq\). Let \(m\in\mathcal{M}\), let \(R_0\) denote the linear-side
randomness, and let \(R_1\in\Rq\) denote the inverse-side randomness. Write
\[
L_B(m,R_0):=B_0+B_1m+B_2R_0.
\]
We say that \((B,m,R_0,R_1)\) is \emph{admissible} if \(B_3-R_1\in \Rq^\times\), equivalently, if
\(B_3-R_1\) is non-zero in every CRT slot of \(\Rq\).

When this admissibility condition holds, the rational hash is the partial function
\[
T=h_{\mathrm{tran},m,(R_0,R_1)}(B)
\iff
\exists Z\in\Rq:\ T=L_B(m,R_0)+Z\ \land\ (B_3-R_1)\Had Z=1.
\]
Equivalently, the eliminated relation is
\begin{equation}\label{eq:hash-cleared}
  (B_3-R_1)\Had\bigl(T-L_B(m,R_0)\bigr)=1.
\end{equation}
\end{definition}

\begin{lemma}[Eliminated form under admissibility]\label{lem:hash-cleared-admissible}
Let
$\Rq \cong \prod_{i=1}^s F_i$
be the CRT decomposition of \(\Rq\) into field factors, and write
\[
L_B(m,R_0)=(\ell_1,\ldots,\ell_s),\qquad
B_3-R_1=(d_1,\ldots,d_s),\qquad
T=(t_1,\ldots,t_s).
\]
If \((B,m,R_0,R_1)\) is admissible, that is, \(d_i\neq 0\) for every \(i\), then
\[
T=h_{\mathrm{tran},m,(R_0,R_1)}(B)
\iff
(B_3-R_1)\Had\bigl(T-L_B(m,R_0)\bigr)=1.
\]
\end{lemma}

\begin{proof}
Under the CRT isomorphism, the defining relation is the coordinate-wise statement
\[
t_i=\ell_i+d_i^{-1}
\qquad\text{for every }i.
\]
If \(d_i\neq 0\) for every \(i\), then each \(d_i\) is invertible in \(F_i\), so
\[
t_i=\ell_i+d_i^{-1}
\iff
d_i(t_i-\ell_i)=1.
\]
Collecting the coordinates yields the equivalence. If admissibility fails, then \(d_i=0\) in
at least one slot, so the quotient relation is undefined there and the eliminated relation is outside its intended
scope.
\end{proof}

Note that the undefined case \(h_{\mathrm{tran},m,(R_0,R_1)}(B)=\bot\) is exactly the inadmissibility event
that \(B_3-R_1\) is not invertible in \(R_q\), and is modeled by the admissibility parameter \(\delta\).

\subsubsection{Aligned commitment and aligned BB-tran target}
\label{subsec:aligned-hiding}
For the blind-issuance surface used later, we parse the full signed message as
\[
\vecmu=(\vecm\Vert \veck)\in \Rq^{\ell_\mu},
\qquad
\ell_\mu=\ell_m+\ell_k,
\qquad
\mA_\mu=[\mA_m\mid \mA_k].
\]
Let the linear-side randomness be split as
\[
R_0:=(S\Vert E),
\qquad
S\in\Rq^{\ell_s},\quad E\in\Rq^{\ell_{\ComScheme}},
\qquad
R_0\gets \mathcal D_0:=\chi_s^{\ell_s}\times \chi_e^{\ell_{\ComScheme}}.
\]
Let the inverse-side randomness be sampled from an admissible distribution
\[
R_1\gets \mathcal D_1^{\mathrm{adm}}
\]
supported on values satisfying \(B_3-R_1\in \Rq^\times\). Equivalently, one may sample from an unconditional
source \(\mathcal D_1\) and restart until admissible, incurring an inadmissibility term
\[
\delta_{\mathrm{adm}}:=\Pr_{R_1\gets \mathcal D_1}[B_3-R_1\notin \Rq^\times].
\]
Define the aligned commitment by
\begin{equation}\label{eq:aligned-commitment}
  c
  := \Commit(\vecmu;R_0)
  := \mA_m\vecm + \mA_k\veck + \mA_s S + E.
\end{equation}
Let \(U\in \Rq^{1\times \ell_{\ComScheme}}\) be a public linear map from the commitment space to \(\Rq\), and
define the BB-tran blocks by
\begin{equation}\label{eq:aligned-blocks}
B_{1,m}:=U\mA_m,
\qquad
B_{1,k}:=U\mA_k,
\qquad
B_{2,s}:=U\mA_s,
\qquad
B_{2,e}:=U.
\end{equation}
For later sections, we use the shorthand
\[
B_1(M+K):=B_{1,m}M+B_{1,k}K,
\qquad
B_2R_0:=B_{2,s}S+B_{2,e}E.
\]
Then the aligned BB-tran linear part is
\begin{equation}\label{eq:aligned-linear-part}
L_B(\vecm,\veck,S,E)
:= B_0 + B_{1,m}\vecm + B_{1,k}\veck + B_{2,s}S + B_{2,e}E
= B_0 + Uc.
\end{equation}
Define the inverse witness
\begin{equation}\label{eq:aligned-inverse}
Z:=(B_3-R_1)^{-1}\in \Rq,
\qquad\text{so that}\qquad
(B_3-R_1)\Had Z = 1.
\end{equation}
Finally define the aligned target
\begin{equation}\label{eq:aligned-target}
T
:= h_{\mathrm{tran},\vecmu,(R_0,R_1)}(B)
:= L_B(\vecm,\veck,S,E)+Z.
\end{equation}
Equivalently,
\begin{equation}\label{eq:aligned-target-public}
T = B_0 + B_{1,m}\vecm + B_{1,k}\veck + B_{2,s}S + B_{2,e}E + Z = B_0 + Uc + Z.
\end{equation}

\paragraph{Hiding-view experiment.}
For a PPT adversary \(\mathcal A=(\mathcal A_1,\mathcal A_2)\), define
\(\mathsf{Exp}^{\mathrm{hide\mbox{-}view}}_{h_{\mathrm{tran}},\mathcal A}(\lambda)\) as follows.
\begin{enumerate}[leftmargin=1.25em]
  \item Sample \(\mA_m,\mA_k,\mA_s,U,B_0,B_3\) as above and give
  \[
  pp_{\mathrm{hid}} := (\mA_m,\mA_k,\mA_s,U,B_0,B_3)
  \]
  to \(\mathcal A_1\).
  \item \(\mathcal A_1(pp_{\mathrm{hid}})\) outputs two challenge messages
  \[
  \vecmu_0=(\vecm_0\Vert \veck_0),\qquad \vecmu_1=(\vecm_1\Vert \veck_1),
  \]
  together with state \(\mathsf{st}\).
  \item The challenger samples
  \[
  b\getsunif\{0,1\},\qquad (S,E)\gets \mathcal D_0,\qquad R_1\gets \mathcal D_1^{\mathrm{adm}}.
  \]
  \item It computes
  \[
  c^\star := \mA_m\vecm_b + \mA_k\veck_b + \mA_s S + E,
  \]
  \[
  Z^\star := (B_3-R_1)^{-1},
  \]
  and
  \[
  T^\star := B_0 + Uc^\star + Z^\star.
  \]
  \item It gives \((c^\star,T^\star)\) to \(\mathcal A_2\), which outputs a bit \(b'\).
\end{enumerate}
The experiment returns \(1\) iff \(b'=b\). Define
\begin{equation}\label{eq:hide-view-adv}
\Adv^{\mathrm{hide\mbox{-}view}}_{h_{\mathrm{tran}}}(\mathcal A)
:=
\left|
\Pr\!\left[\mathsf{Exp}^{\mathrm{hide\mbox{-}view}}_{h_{\mathrm{tran}},\mathcal A}(\lambda)=1\right]
-\frac12
\right|.
\end{equation}

\paragraph{Target-only hiding experiment.}
Define \(\mathsf{Exp}^{\mathrm{hide\mbox{-}target}}_{h_{\mathrm{tran}},\mathcal A}(\lambda)\) identically,
except that the challenger gives only \(T^\star\) to \(\mathcal A_2\). Let
\begin{equation}\label{eq:hide-target-adv}
\Adv^{\mathrm{hide\mbox{-}target}}_{h_{\mathrm{tran}}}(\mathcal A)
\end{equation}
denote the corresponding advantage.

\begin{theorem}[Commitment-lifted hiding of aligned BB-tran]\label{thm:tran-commit-hiding}
Let \(\Commit\) be the aligned commitment from \eqref{eq:aligned-commitment}. For every PPT adversary
\(\mathcal A\) against \(\mathsf{Exp}^{\mathrm{hide\mbox{-}view}}_{h_{\mathrm{tran}},\mathcal A}\), there exists a PPT
adversary \(\mathcal B\) against hiding of \(\Commit\) such that
\begin{equation}\label{eq:tran-commit-hiding-red}
\Adv^{\mathrm{hide\mbox{-}view}}_{h_{\mathrm{tran}}}(\mathcal A)
=
\Adv^{\mathrm{hide}}_{\ComScheme}(\mathcal B).
\end{equation}
Consequently, if \(\Commit\) is hiding, then the aligned BB-tran issuer view \((c,T)\) is hiding.
\end{theorem}

\begin{proof}
Construct \(\mathcal B=(\mathcal B_1,\mathcal B_2)\) for the commitment-hiding game as follows.

\(\mathcal B_1\) samples
\[
U\getsunif \Rq^{1\times \ell_{\ComScheme}},\qquad B_0,B_3\getsunif \Rq,
\]
and receives from the commitment challenger the public commitment parameters
\[
pp_{\ComScheme}=(\mA_m,\mA_k,\mA_s).
\]
It forwards
\[
pp_{\mathrm{hid}}=(\mA_m,\mA_k,\mA_s,U,B_0,B_3)
\]
to \(\mathcal A_1\). When \(\mathcal A_1\) outputs
\[
\vecmu_0=(\vecm_0\Vert \veck_0),\qquad \vecmu_1=(\vecm_1\Vert \veck_1),\qquad \mathsf{st},
\]
\(\mathcal B_1\) forwards the same message pair to its commitment-hiding challenger.

Let \(c^\star\) be the commitment challenge returned by that challenger, so by definition there is a hidden bit \(b\)
such that
\[
c^\star = \mA_m\vecm_b + \mA_k\veck_b + \mA_s S + E
\]
for \((S,E)\gets \mathcal D_0\).

Now \(\mathcal B_2\) samples
\[
R_1\gets \mathcal D_1^{\mathrm{adm}},
\qquad
Z^\star := (B_3-R_1)^{-1},
\]
and computes
\begin{equation}\label{eq:tran-commit-red-target}
T^\star := B_0 + Uc^\star + Z^\star.
\end{equation}
It gives \((c^\star,T^\star)\) to \(\mathcal A_2\) and outputs whatever bit \(b'\) \(\mathcal A_2\)
returns.

It remains to check that the distribution seen by \(\mathcal A\) is exactly the real one. Using
\eqref{eq:aligned-blocks}, \eqref{eq:aligned-linear-part}, and the defining equation for \(c^\star\),
\begin{equation}\label{eq:tran-commit-red-distribution}
\begin{aligned}
T^\star
&= B_0 + U(\mA_m\vecm_b + \mA_k\veck_b + \mA_s S + E) + Z^\star \\
&= B_0 + (U\mA_m)\vecm_b + (U\mA_k)\veck_b + (U\mA_s)S + UE + Z^\star \\
&= B_0 + B_{1,m}\vecm_b + B_{1,k}\veck_b + B_{2,s}S + B_{2,e}E + Z^\star \\
&= h_{\mathrm{tran},\vecmu_b,((S,E),R_1)}(B).
\end{aligned}
\end{equation}
So the pair \((c^\star,T^\star)\) is distributed exactly as in
\(\mathsf{Exp}^{\mathrm{hide\mbox{-}view}}_{h_{\mathrm{tran}},\mathcal A}\). Therefore the advantages are equal,
proving \eqref{eq:tran-commit-hiding-red}.
\end{proof}

\begin{theorem}[LWE-based hiding of aligned BB-tran]\label{thm:tran-lwe-hiding}
Assume the decisional module-LWE problem for the distribution pair \((\chi_s,\chi_e)\) is hard in the
following form: given
\[
\mA_s\getsunif \Rq^{\ell_{\ComScheme}\times \ell_s}
\]
and a challenge vector \(Y\in\Rq^{\ell_{\ComScheme}}\), distinguish between
\begin{equation}\label{eq:tran-lwe-real}
Y = \mA_s S + E
\quad\text{with}\quad
S\gets \chi_s^{\ell_s},\ E\gets \chi_e^{\ell_{\ComScheme}},
\end{equation}
and
\begin{equation}\label{eq:tran-lwe-rand}
Y \getsunif \Rq^{\ell_{\ComScheme}}.
\end{equation}
Then for every PPT adversary \(\mathcal A\),
\begin{equation}\label{eq:tran-lwe-adv}
\Adv^{\mathrm{hide\mbox{-}view}}_{h_{\mathrm{tran}}}(\mathcal A)
\le
\Adv^{\mathrm{d\mbox{-}MLWE}}_{q,\ell_s,\ell_{\ComScheme},\chi_s,\chi_e}(\mathcal C)
\end{equation}
for some PPT distinguisher \(\mathcal C\). In fact the reduction below is exact.
\end{theorem}

\begin{proof}
Build a distinguisher \(\mathcal C\) for decisional module-LWE. It receives a challenge pair
\((\mA_s,Y)\), where \(Y\) is either an LWE sample as in \eqref{eq:tran-lwe-real} or uniform as in
\eqref{eq:tran-lwe-rand}.

First, \(\mathcal C\) samples independently
\[
\mA_m\getsunif \Rq^{\ell_{\ComScheme}\times \ell_m},\qquad
\mA_k\getsunif \Rq^{\ell_{\ComScheme}\times \ell_k},
\]
\[
U\getsunif \Rq^{1\times \ell_{\ComScheme}},\qquad B_0,B_3\getsunif \Rq.
\]
It gives
\[
pp_{\mathrm{hid}}=(\mA_m,\mA_k,\mA_s,U,B_0,B_3)
\]
to \(\mathcal A_1\), which returns challenge messages
\[
\vecmu_0=(\vecm_0\Vert \veck_0),\qquad \vecmu_1=(\vecm_1\Vert \veck_1),
\]
and state \(\mathsf{st}\).

Then \(\mathcal C\) samples
\[
b\getsunif\{0,1\},\qquad R_1\gets \mathcal D_1^{\mathrm{adm}},\qquad Z^\star := (B_3-R_1)^{-1},
\]
and sets
\begin{equation}\label{eq:tran-lwe-commitment}
c^\star := \mA_m\vecm_b + \mA_k\veck_b + Y,
\end{equation}
\begin{equation}\label{eq:tran-lwe-target}
T^\star := B_0 + Uc^\star + Z^\star.
\end{equation}
It gives \((c^\star,T^\star)\) to \(\mathcal A_2\), receives a guess \(b'\), and outputs \(1\) iff
\(b'=b\).

If \(Y=\mA_sS+E\) is an LWE sample, then
\[
c^\star = \mA_m\vecm_b + \mA_k\veck_b + \mA_sS + E,
\]
hence by the same calculation as in \eqref{eq:tran-commit-red-distribution},
\[
T^\star = h_{\mathrm{tran},\vecmu_b,((S,E),R_1)}(B),
\]
so \((c^\star,T^\star)\) is exactly the real hiding-view challenge. Therefore
\begin{equation}\label{eq:tran-lwe-success-real}
\Pr[\mathcal C \text{ outputs }1 \mid Y=\mA_sS+E]
=
\frac12 + \Adv^{\mathrm{hide\mbox{-}view}}_{h_{\mathrm{tran}}}(\mathcal A).
\end{equation}
If instead \(Y\getsunif \Rq^{\ell_{\ComScheme}}\), then for each fixed \(b\),
\[
c^\star = \mA_m\vecm_b + \mA_k\veck_b + Y
\]
is still uniform on \(\Rq^{\ell_{\ComScheme}}\), because translating a uniform distribution by a fixed public
offset preserves uniformity. Hence \(c^\star\) is independent of \(b\). Since \(R_1\) is sampled independently of
\(b\), so are \(Z^\star\) and
\[
T^\star = B_0 + Uc^\star + Z^\star.
\]
Thus the whole pair \((c^\star,T^\star)\) is independent of the hidden bit \(b\), and \(\mathcal A_2\) has
success probability exactly \(1/2\):
\begin{equation}\label{eq:tran-lwe-success-rand}
\Pr[\mathcal C \text{ outputs }1 \mid Y\getsunif \Rq^{\ell_{\ComScheme}}]
=
\frac12.
\end{equation}
Subtracting \eqref{eq:tran-lwe-success-rand} from \eqref{eq:tran-lwe-success-real}, the distinguishing
advantage of \(\mathcal C\) is exactly \eqref{eq:tran-lwe-adv}.
\end{proof}

\begin{corollary}[The rational-hash output itself is hiding]\label{cor:tran-target-hiding}
For every PPT adversary \(\mathcal A\) in the target-only game,
\begin{equation}\label{eq:tran-target-hiding}
\Adv^{\mathrm{hide\mbox{-}target}}_{h_{\mathrm{tran}}}(\mathcal A)
\le
\Adv^{\mathrm{hide\mbox{-}view}}_{h_{\mathrm{tran}}}(\mathcal A')
\le
\Adv^{\mathrm{d\mbox{-}MLWE}}(\mathcal C)
\end{equation}
for suitable PPT adversaries \(\mathcal A',\mathcal C\).
\end{corollary}

\begin{proof}
Given a target-only adversary \(\mathcal A\), build a view adversary \(\mathcal A'\) that simply ignores the
commitment component \(c^\star\) and passes only \(T^\star\) to \(\mathcal A\). This preserves the success probability
exactly. Then apply Theorem~\ref{thm:tran-lwe-hiding}.
\end{proof}

\begin{remark}\label{rem:aligned-public-identity}
The privacy statement established for issuance is issuer-view hiding of the public issuance surface. Full
blindness in the two-session sense of Definition~\ref{def:bs-blindness} is left as a separate problem.

The load-bearing public identity is
\[
T = B_0 + Uc + (B_3-R_1)^{-1}.
\]
The issuance syntax in Section~\ref{subsec:bs-protocol} is written so that this identity holds literally.
\end{remark}

We now  define a signature scheme for a Trapdoor Rational Hash.
%\cite[Figure~4 and Section~4.1]{EPRINT:DKLW25} in 

\begin{definition}[vSIS-style hash-and-sign]\label{def:vsis-template}
Let \(\TDRH\) be a trapdoor rational hash scheme  with parameters
\(\PublicParamsTrapdoor\). The associated signature scheme
\(\vsisscheme=(\KeyGen,\Sign,\SigVerify)\) is defined by
\begin{itemize}[leftmargin=1.25em]
  \item \(\KeyGen(\SecParam)\) runs \((\mA,\trapdoor)\gets \TrapGen(\PublicParamsTrapdoor)\) and outputs
  \(\vk=\mA\), \(\sk=\trapdoor\).
  \item \(\Sign(\PublicParamsTrapdoor,\sk,m)\) for \(m\in\mathcal{M}_B\) samples \(\chi\gets\mathcal{X}\) until
  \(\vect=\Eval(B,m,\chi)\neq\bot\), computes \(\vecu\gets \SampPre(\trapdoor,\vect)\), and outputs
  \(\sigma=(\vecu,\chi)\).
  \item \(\SigVerify(\PublicParamsTrapdoor,\vk,m,\sigma)\) for \(\sigma=(\vecu,\chi)\) accepts iff \(m\in\mathcal{M}_B\),
  \(\chi\in\mathcal{X}\), \(\Eval(B,m,\chi)\neq\bot\), \(\|\vecu\|_\infty\le \BoundSig\), and
  \[
    \mA\vecu=\Eval(B,m,\chi).
  \]
\end{itemize}
\end{definition}

Correctness requires that all admissible parameters $\PublicParamsTrapdoor$ generated by $\Setup()$, correctly generated $(\vk,\sk)$ and \(m\in\mathcal{M}_B\), the signatures generated using $\Sign$ are accepted by $\SigVerify$ except with
negligible probability.
\paragraph{Security of vSIS-style signatures.}




\begin{comment}
\begin{definition}[Hinted-vSIS and \(\mathsf{GenISIS}_f\) variants from DKLW25]\label{def:dklw-assumptions}
We recall the DKLW25 assumption variants used in the reduction discussion below.
Fix a public left factor \(F\), hint space \(H\), target space \(G\), predicate \(P\), and query bound \(Q\), all as in
\cite[Definition~5]{EPRINT:DKLW25}.
\begin{itemize}[leftmargin=1.25em]
  \item \textbf{\(\mathsf{Hint\mbox{-}vSIS}\).} The adversary chooses a query set \(\mathcal{Q}\subseteq H\) and a fresh
  target function \(g^\star\in G\setminus \mathcal{Q}\) before seeing the sampled matrix \(A\). It then receives short
  preimages of the queried hints \(q(A)\) under \(F(A)\) for every \(q\in\mathcal{Q}\), and must output a non-zero short
  vector \(u^\star\) such that \(F(A)u^\star=g^\star(A)\). The instance is restricted by the predicate condition
  \(P(F\cup \mathcal{Q}\cup(\{g^\star\}\setminus\{0\}))=1\).
  \item \textbf{\(s\)-\(\mathsf{Hint\mbox{-}vSIS}\).} This is the strong variant of the same experiment. The adversary still
  chooses the hint set \(\mathcal{Q}\) before seeing \(A\), but it is allowed to choose the fresh target \(g^\star\) only
  after seeing \(A\) and the hinted preimages.
  \item \textbf{\(s\)-\(\$\)-\(\mathsf{Hint\mbox{-}vSIS}\).} This is the strong random-hinted variant. It is the same as
  \(s\)-\(\mathsf{Hint\mbox{-}vSIS}\), except that the query set \(\mathcal{Q}\) is sampled uniformly at random from \(H\)
  rather than chosen by the adversary \cite[Definition~5]{EPRINT:DKLW25}.
\end{itemize}

Now fix a keyed evaluation family \(f:K\times M\times X\to R_q^n\) as in \cite[Definition~8]{EPRINT:DKLW25}.
\begin{itemize}[leftmargin=1.25em]
  \item \textbf{\(\mathsf{GenISIS}_f\).} The challenger samples a public key \(\kappa\gets K\) and a uniform matrix
  \(A\getsunif R_q^{n\times m}\). It provides a preimage oracle that, on input \(\mu\in M\), samples \(\chi\gets X\),
  sets \(t=f(\kappa,\mu,\chi)\), returns a short preimage \(s\) such that \(As=t\), and records the tuple
  \((s,\mu,\chi)\). The adversary wins if it outputs a fresh tuple \((s^\star,\mu^\star,\chi^\star)\) not previously
  returned by the oracle such that \(As^\star=f(\kappa,\mu^\star,\chi^\star)\) and \(0<\|s^\star\|\le \beta\).
  \item \textbf{\(\mathsf{IntGenISIS}_f\).} DKLW25 also defines an interactive extension in which the challenger offers a
  signing-style oracle with an additional linear masking term, and Theorem~6 shows how this interactive problem reduces
  back to \(\mathsf{GenISIS}_f\) under enlarged parameters.
\end{itemize}
\end{definition}
\end{comment}
%We use the signature-security experiments of \cite[Definition~4 and Figure~3]{EPRINT:DKLW25} specialized to  Definition~\ref{def:vsis-template} and 

We now define security as strong existential unforgeability under random-message attacks.
\begin{definition}[Unforgeability for vSIS-style signatures]\label{def:vsis-games}
Let \(\TDRH\) be a trapdoor rational hash scheme with parameters
\(\PublicParamsTrapdoor\) and
\(\vsisscheme=(\KeyGen,\Sign,\SigVerify)\) be the associated signature scheme.
Fix \(Q\in\poly(\lambda)\). In the experiment
\(\mathsf{Exp}_{\mathrm{sEUF\mbox{-}RMA}}^{\Sigma_{\mathrm{vSIS}},\adversary,Q}(\lambda)\), the challenger runs
\((\vk,\sk)\gets \KeyGen\), then uniformly samples \(m_1,\ldots,m_Q \gets \mathcal{M}_B\) and computes
\(\sigma_i\leftarrow \Sign(\sk,m_i)\) for each \(i\in [Q]\). Finally
\(\adversary\) on input $\vk,(m_i,\sigma_i)_{i\in [Q]}$ and the trapdoor rational hash parameters outputs a pair \((m^\star,\sigma^\star)\) and wins iff:
\begin{enumerate}[leftmargin=1.25em]
    \item \(\SigVerify(\vk,m^\star,\sigma^\star)=1\), and
    \item \((m^\star,\sigma^\star)\notin \{(m_i,\sigma_i):1\le i\le Q\}\).
\end{enumerate}
\end{definition}
We can equivalently define \(\mathrm{sEUF\mbox{-}CMA}\), \(\mathrm{sEUF\mbox{-}SMA}\), and
\(\mathrm{sSUF\mbox{-}SMA}\) security where \(\mathrm{CMA}\)-security allows the adversary to choose all queried
messages after seeing \(\vk\), while \(\mathrm{SMA}\)-security requires the adversary to fix them before seeing
\(\vk\).

\begin{lemma}[Imported security properties of the non-blind rational-hash signature layer]\label{lem:vsis-unforgeability}
Let \(\ell_{\mathrm{hash}}=4\) for the present ARC instantiation and set
\[
  k_{\mathrm{DKLW}}:=m+\ell_{\mathrm{hash}}.
\]
Assume the following:
\begin{enumerate}[leftmargin=1.25em]
    \item \(\PublicParamsTrapdoor\) are admissible parameters for \((\Setup,\TrapGen,\SampPre)\);
    \item the inadmissibility bound satisfies \(\delta\le \negl(\lambda)\);
    \item a DKLW-compatible norm bound \(\beta_{\mathrm{DKLW}}\) is fixed so that every signature accepted by
    the present verifier also satisfies \(\|\vecu\|\le \beta_{\mathrm{DKLW}}\) in the norm used in
    \cite{EPRINT:DKLW25};
    \item the predicate condition of \cite[Definition~7]{EPRINT:DKLW25} holds for the instantiated rational-hash
    family; and
    \item the denominator-bound, strong-linear-independence, and adaptive-security side conditions required by
    \cite[Theorems~4--8, Corollary~1, and Section~5.2]{EPRINT:DKLW25} hold for the instantiated family.
\end{enumerate}
Here \(\beta_{\mathrm{DKLW}}\) is the verifier norm parameter from the imported DKLW theorems. It need not literally
equal \(\BoundSig\), but it must dominate every signature accepted by our verifier when measured in the norm convention
of \cite{EPRINT:DKLW25}. Assumptions~(4) and~(5) are precisely the instantiated BB-tran compatibility conditions needed
to invoke the cited DKLW reduction chain.
Then the following security consequences hold.
\begin{enumerate}[leftmargin=1.25em]
    \item The non-blind template of Definition~\ref{def:vsis-template} is correct.
    \item Under the strong random-hinted vSIS assumption for the induced parameter tuple, the non-blind
    signature scheme is \(\mathrm{sEUF\mbox{-}RMA}\)-secure.
    \item For \(f(B,m,\chi)=h_{\mathrm{tran},m,\chi}(B)\), the experiment \(\mathsf{GenISIS}_f\) is
    equivalent to the \(\mathrm{sEUF\mbox{-}RMA}\) experiment.
    \item For the BB-tran-type instantiation used here, the \(\mathsf{GenISIS}_f\) /
    \(\mathsf{IntGenISIS}_f\) framework of \cite[Section~5.2 and Theorem~6]{EPRINT:DKLW25} gives the corresponding
    adaptive/CMA security statement for the non-blind signature layer under the stated parameter conditions.
    \item Under the strong-hinted vSIS assumption, the same non-blind family also enjoys
    \(\mathrm{sEUF\mbox{-}SMA}\) security.
\end{enumerate}
\end{lemma}

\begin{proof}
Items~(1), (2), (3), and~(5) follow from
\cite[Theorems~4, 5, 7, 8, and Corollary~1]{EPRINT:DKLW25} under the stated specialization. Item~(4) is the
BB-tran-specific adaptive-security lift of \cite[Section~5.2 and Theorem~6]{EPRINT:DKLW25}, which places the
non-blind signature layer on the \(\mathsf{GenISIS}_f\)/\(\mathsf{IntGenISIS}_f\) interface used there for
\(\mathrm{sEUF\mbox{-}CMA}\) security.
\end{proof}

\begin{remark}[Scope of the DKLW import]\label{rem:vsis-unforgeability}
Lemma~\ref{lem:vsis-unforgeability} is the external reduction basis for the non-blind signature layer. It
does not by itself imply blindness or one-more unforgeability of the blind issuance protocol, and it does not by
itself yield an ARC security theorem.
\end{remark}

\subsection{Non-interactive zero-knowledge arguments of knowledge (NIZKs)}
\label{subsec:prelim-nizk}
\label{subsec:nizk-ark}
We use the NIZK-ARK syntax of the SmallWood framework instantiated in the random-oracle model
\cite[Definition~3]{EPRINT:FenRiv25b}.

\begin{definition}[NIZKs in the random oracle model]\label{def:nizk-ark}
Let $\mathcal{R}\subseteq \mathcal{X}\times \mathcal{W}$  be a family of relations.
A (transparent) non-interactive zero-knowledge \emph{argument of knowledge} for $\mathcal{R}$ in the random oracle model
(with oracle $\RO(\cdot)$) is a pair of algorithms $(\ZKProve^{\RO},\ZKVerify^{\RO})$:
\begin{itemize}[leftmargin=1.25em]
  \item $\ZKProve^{\RO}(w ,x)\to \pi$ is probabilistic and on input $(x,w)$ with
  $(x,w)\in R$ outputs a proof $\pi$.
  \item $\ZKVerify^{\RO}(\pi,x)\to \{0,1\}$ is a deterministic polynomial-time algorithm.
\end{itemize}

\paragraph{Completeness.}
For all $(x,w)\in\mathcal{R}$,
\[
  \Pr\big[\ZKVerify^{\RO}(\pi,x)=1 \ :\ \pi\gets \ZKProve^{\RO}(w,x)\big]=1.
\]

\paragraph{Straightline-extractable knowledge soundness.}
There exists a PPT algorithm $\Ext$, also called \emph{extractor}, such that for any PPT adversary $\Adv$,
\[
  \Pr\Big[
    \ZKVerify^{\RO}(\pi,x)=1 \ \wedge\ (x,w)\notin R
  \Big]\ \le\ \varepsilon,
\]
where $(\pi,Q)\leftarrow \Adv^{\RO}(x)$, $Q$ is the set of query--response pairs made to $\RO$ by $\Adv$,
and $w\leftarrow \Ext(Q,\pi)$.

\paragraph{Zero knowledge.}
There exists a PPT algorithm $\Sim$, called the \emph{simulator}, such that for any $(x,w)\in \mathcal{R}$ and PPT algorithm $\Adv$:
\[
  \Big|\Pr[\Adv^{\RO}(\pi_{\mathsf{real}})=1]-\Pr[\Adv^{\RO}(\pi_{\mathsf{sim}})=1]\Big|
  \le \xi,
\]
where $\pi_{\mathsf{real}}\gets \ZKProve^{\RO}(w ,x)$,  $\pi_{\mathsf{sim}}\gets \Sim^{\RO}(x)$ and $\Sim$ can program the output of $\RO$ on any input.
\end{definition}

We instantiate these proofs using the SmallWood PACS/PIOP framework compiled with Fiat--Shamir in the random oracle
model \cite{EPRINT:FenRiv25b}. SmallWood is a commit-and-prove proof system that only relies on symmetric primitives. While its proof size is linear in $|w|$ it is sublinear in the circuit size $x$.

%At a high level, SmallWood expresses statements as polynomial constraints evaluated
%on an $\Omega\subset\Fq$ packing domain, supports both \emph{parallel} constraints checked point-wise on $\Omega$ and
%\emph{aggregated} constraints checked via an $\Omega$-sum identity, and uses random verifier queries outside $\Omega$ to
%spot-check correctness of committed witness polynomials while preserving HVZK via blinding evaluations
%(Section~\ref{sec:smallwood-model}). 



\subsection{Blind signatures}
A blind signature algorithm allows a holder $\Holder$ to obtain signatures valid under the verification key of an issuer $\Issuer$, without the latter learning the signed message $\mu$. Our definitions follow \cite{C:Fischlin06}. When any algorithm $\mathcal{B}$ uses the interactive protocol $\Issue$ with $\Issuer$, we will write $\mathcal{B}^{\Issue:\langle \Issuer(\cdot):\cdot \rangle}$ and adjust this the notation for other interactions accordingly.

\begin{definition}[Blind Signature]\label{def:bs-syntax}
	A Blind Signature scheme is a tuple of algorithms
	\[
	\BlindSignature=(\Setup,\IKeyGen,\Issue,\Verify)
	\]
	such that:
	\begin{itemize}[leftmargin=1.25em]
		\item $\Setup(\SecParam)$ is a PPT algorithm that outputs public parameters $\PublicParamsBS$. All other algorithms obtain $\SecParam,\PublicParamsBS$ as implicit inputs.
		\item $\IKeyGen()$ is a PPT algorithm that outputs issuer keys $(\vk,\sk)$.
		\item $\Issue$ is an interactive protocol between $\Issuer$ and $\Holder$, where $\Issuer$ has input $\sk$ and $\Holder$ has input $\mu,\vk$. At the end of the protocol, $\Holder$ obtains $\sig$.
		\item $\Verify(\vk,\mu,\sig)$ is a deterministic polynomial-time algorithm that outputs a bit.
	\end{itemize}
\end{definition}

\begin{definition}[Correctness]\label{def:bs-correctness}
	A blind signature scheme is correct if there exists \(\negl\) such that for all messages \(\mu\),
	\[
	\Pr\!\left[
	\begin{array}{l}
		\PublicParamsBS \gets \Setup(\SecParam); \\
		(\vk,\sk)\gets \IKeyGen(\PublicParamsBS);\\
		\sig\gets \Holder^{\Issue: \langle \Issuer(\sk): \cdot  \rangle }(\mu,\vk);\\
		\Verify(\vk,\mu,\sig)=1
	\end{array}
	\right]
	\ge 1-\negl(\lambda).
	\]
\end{definition}


\begin{definition}[Blindness]\label{def:bs-blindness}
  Let $\mathsf{BSig}$ be a blind signature scheme and let $\mu_0 \neq \mu_1$ be any two messages. Define the
  experiment $\mathsf{Exp}^{\mathsf{sig\mbox{-}blind}}_{\mathsf{BSig},\Adv}(\lambda)$ between a challenger
  $\Challenger$ and a three-stage stateful adversary $\Adv=(\Adv_1,\Adv_2,\Adv_3)$ as follows.
 \begin{enumerate}[leftmargin=1.25em]
   \item $\Challenger$ runs $\PublicParamsBS\gets \Setup(\SecParam)$.
   \item $\Adv_1(\PublicParamsBS)$ outputs a verification key $\vk$.
	   \item Run
			$\Adv_2^{\Issue:\langle \cdot: \Holder(\mu_0,\vk)\rangle ,~ \Issue:\langle
	    \cdot: \Holder(\mu_1,\vk)\rangle };$
	     at the end of which $\Challenger$ obtains $\sig_0,\sig_1$.
	   \item If $\bot\in\{\sig_0,\sig_1\}$ then output a uniform bit and halt.
	   \item Otherwise, $\Challenger$ samples $b\getsunif\{0,1\}$ and gives the ordered pair
	   \((\sig_b,\sig_{1-b})\) to $\Adv_3$, which
	   outputs a bit $b'$.
	 \end{enumerate}
 Define $\Adv^{\mathsf{sig\mbox{-}blind}}_{\mathsf{BSig}}(\Adv) := \big|\Pr[b'=b]-\tfrac12\big|$.
 We say $\mathsf{BSig}$ is \emph{blind} if for all PPT $\Adv$ this advantage is negligible in $\lambda$.
\end{definition}



\begin{definition}[One-more unforgeability]\label{def:bs-euf-cma}
A blind signature scheme is one-more unforgeable if there exists \(\negl\) such that for any PPT adversary
\(\adversary^{\Issue: \langle \Issuer(\sk): \cdot \rangle}(\vk)\) that makes at most \(Q=Q(\lambda)\) oracle queries to \(\Issue\),
\[
  \Pr\!\left[
    \begin{array}{l}
     	\PublicParamsBS \gets \Setup(\SecParam); \\
      (\vk,\sk)\gets \IKeyGen(\PublicParamsBS);\\
      \{(\mu_i,\sig_i)\}_{i\in[Q+1]}\gets \adversary^{\Issue \langle \Issuer(\sk):\cdot \rangle}(\vk):\\
      \forall\,1\le i<j\le Q+1:\ \mu_i\neq \mu_j\ \land\ \\ \forall\,i\in[Q+1]:\ \Verify(\vk,\mu_i,\sig_i)=1
    \end{array}
  \right]
  \le \negl(\lambda).
\]
\end{definition}

%Our issuance protocol (Section~\ref{sec:blind-signature}) instantiates this primitive and additionally packages the
%interaction with pre-sign and post-sign NIZKs tailored to ARC.

\paragraph{Rate limiting at the ARC layer.}
We do not introduce a separate stand-alone rate-limiting primitive. In the present paper, rate limiting is part of the
ARC syntax itself: the holder presents a credential together with a public nonce and a PRF-derived public tag, and the
verifier enforces the policy through local state on previously accepted \((\nonce,\prftag)\) pairs.
\subsection{Anonymous Rate-Limiting Credentials (ARCs)}
\label{subsec:arc-security}
Anonymous rate-limited credentials~\cite{ietf-privacypass-arc-protocol-01} extend blind signatures by adding a showing algorithm \(\Show\) that may create
multiple presentations of a previously issued credential. We explicitly introduce the verifier party \(\Verifier\),
which receives and checks those presentations. Our definition follows the established syntax and security models
of~\cite{JC:FucHanSla19,C:BLNS23}, specialized to the ARC setting considered here.

We keep \(\Setup\) and \(\IKeyGen\) separate, so the public commitment matrix is sampled independently of the issuer
keys. We also separate \(\Show\) and \(\Verify\) to distinguish the client-side presentation algorithm from the
verifier-side stateful check.
\begin{definition}[ARC scheme syntax]\label{def:arc-syntax}
Let $\nonceSpace$ be a finite set of nonces.
An anonymous rate-limited credential (ARC) scheme is a tuple of algorithms
\[
  \ARC=(\Setup,\IKeyGen,\Issue,\Show,\Verify)
\]
such that:

\begin{itemize}[leftmargin=1.25em]
   \item $\Setup(\SecParam)$ is a PPT algorithm that outputs public parameters $\PublicParamsARC$. All other algorithms obtain $\SecParam,\PublicParamsARC$ as implicit inputs.

    \item $\IKeyGen()$ is a PPT algorithm that outputs issuer keys $(\vk,\sk)$.
\item $\Issue$ is an interactive protocol between $\Issuer$ and $\Holder$, where $\Issuer$ has input $\sk$ and $\Holder$ has input a message $\mu$ and verification key $\vk$. At the end of the protocol, $\Holder$ obtains the credential $\credential$.
  \item $\Show(\credential,\nonce)$ on input a credential $\credential$ and a nonce $\nonce\in \nonceSpace$ outputs a presentation $\presentation$.

  \item $\Verify(\vk,\presentation,\algstate)$ is a deterministic polynomial-time algorithm that on input a
  verification key $\vk$, a presentation $\presentation$, and a verifier state $\algstate$ outputs a validity bit
  together with an updated state $\algstate'$.
	 \end{itemize}
\end{definition}

Adjusting the blind-signature correctness notion of Definition~\ref{def:bs-correctness} to ARC is immediate.
Privacy at the ARC layer is captured below by presentation unlinkability, while issuance-side privacy continues to be
expressed by Definition~\ref{def:bs-blindness} together with the issuer-view hiding property proved in
Section~\ref{sec:blind-signature}.


We additionally require the following soundness notion, which captures that each issuance can support at most one accepting presentation per
nonce.

\begin{definition}[ARC soundness]\label{def:arc-soundness}
Let \(\ARC\) be an ARC scheme with nonce space \(\nonceSpace\). For \(Q=Q(\lambda)\), set
\[
K:=Q\cdot |\nonceSpace|+1.
\]
Consider the following experiment. The challenger samples
\[
\PublicParamsARC\gets \Setup(\SecParam),\qquad (\vk,\sk)\gets \IKeyGen(\PublicParamsARC),
\]
initializes the verifier state \(\algstate_0=\emptyset \), and gives \((\PublicParamsARC,\vk)\) to a PPT adversary
\(\adversary^{\Issue:\langle \Issuer(\sk): \cdot \rangle}\) that makes at most \(Q\) issuance queries. The adversary then
outputs a sequence \((\presentation_i)_{i\in [K]}\). For \(i=1,\ldots,K\), the challenger computes
\[
(b_i,\algstate_i)\gets \Verify(\vk,\presentation_i,\algstate_{i-1}).
\]
We say that \(\ARC\) is \emph{sound} if for every PPT adversary, the probability that \(b_i=1\) for every
\(i\in[K]\) is negligible in \(\lambda\).
\end{definition}

This definition is the ARC analogue of one-more unforgeability. Since each honestly issued credential can be shown for
all \(|\nonceSpace|\) different nonces, \(Q\) issuance interactions may legitimately yield \(Q\cdot |\nonceSpace|\)
accepting presentations. Producing strictly more than that forces a policy violation for at least one nonce or credential by the
pigeonhole principle.

However, we additionally require that if the same credential \(\credential\) is presented multiple times, but for
different nonces from \(\nonceSpace\), then the presentations should be unlinkable. To formalize this, we let a
challenger create two credentials in interactions where the adversary controls the issuer. Then the challenger either
creates two presentations for the same credential, or for different credentials.

\begin{definition}[Presentation unlinkability]\label{def:pres-unlink}
Let $\ARC$ be an ARC scheme. Consider the following experiment
$\mathsf{Exp}^{\mathsf{pres\mbox{-}unlink}}_{\ARC,\Adv}(\lambda)$ between a challenger $\Challenger$ and a three-stage stateful adversary $\Adv=(\Adv_1,\Adv_2,\Adv_3)$:
\begin{enumerate}[leftmargin=1.25em]
  \item $\Challenger$ runs $\PublicParamsARC\gets \Setup(\SecParam)$.
  \item Run $(\vk,\mu)\gets \Adv_1(\PublicParamsARC)$.
  \item Run $\Adv_2^{\Issue:\langle \cdot: \Challenger(\vk,\mu)\rangle ,~ \Issue:\langle \cdot: \Challenger(\vk,\mu)\rangle }$, at the end of which $\Challenger$ obtains $\credential_0,\credential_1$ and $\Adv_2$ outputs $\nonce_0,\nonce_1$.
  \item If $\credential_0$ or $\credential_1$ is $\bot$ or if $\nonce_0=\nonce_1$ then output a random bit.

  \item $\Challenger$ samples $b\getsunif\{0,1,2\}$ and computes
  \begin{align*}
    \presentation_0 \leftarrow & \Show(\credential_0,\nonce_0) \\
    \presentation_1 \leftarrow & \begin{cases}\Show(\credential_0,\nonce_1) & \text{ if } b=0\\
    \Show(\credential_1,\nonce_0) & \text{ if } b=1\\
    \Show(\credential_1,\nonce_1) & \text{ if } b=2
    \end{cases}
  \end{align*}
  \item $\Adv_3(\presentation_0,\presentation_1)$ outputs $b'$.
\end{enumerate}
Define $\Adv^{\mathsf{pres\mbox{-}unlink}}_{\ARC}(\Adv) := \big| \Pr[b'=b]-\tfrac13 \big|$.
We say ARC has \emph{presentation unlinkability} if for all PPT $\Adv$ this advantage is negligible in $\lambda$.

\end{definition}


%\paragraph{Nonce reuse caveat (intended).}
%If a nonce is reused, deterministic presentation components (e.g., a
%rate-limiting tag) breaks unlinkability for those presentations; therefore the experiment enforces nonce freshness.

\subsection{Pseudorandom functions}
\label{subsec:prelim-prf}
\label{subsec:prf}


\begin{definition}[Pseudorandom function (PRF)]\label{def:prf}
Let $\mathcal{K}$, $\mathcal{D}$ and $\mathcal{T}$ be sets such that $|\mathcal{K}|$ is exponential in $\lambda$.
A function family $\PRF:\mathcal{K}\times\mathcal{D}\to \mathcal{T}$ is a PRF if for all PPT adversaries
$\Adv$,
\[
  \Adv^{\mathsf{prf}}_{\PRF}(\Adv)
  := \Big|\Pr[\Adv^{\mathcal{O}_0}=1]-\Pr[\Adv^{\mathcal{O}_1}=1]\Big|
  \le \negl(\lambda),
\]
where $\mathcal{O}_0(\cdot)=\PRF(k,\cdot)$ for uniform $k\getsunif \mathcal{K}$ and $\mathcal{O}_1(\cdot)$ is a uniform
random function from $\mathcal{D}$ to $\mathcal{T}$.
\end{definition}

%\paragraph{ARC-SPRUCE instantiation rule.}
%In ARC-SPRUCE, the PRF key is the secret message component \(k\) inside the signed message \(m=\mu\Vert k\), and the
%PRF input includes the public nonce (and, if desired, a policy context encoding).

In this work, we instantiate PRFs based on the Poseidon2 construction~\cite{AFRICACRYPT:GraKhoSch23}.

It is based on the Poseidon2 permutation, which is defined as follows:
\begin{definition}[Poseidon2-style permutation structure]\label{def:poseidon-perm}
		Let $q$ be a prime, $d\ge 3$ such that $\gcd(d,q-1)=1$ and $t,R_F,R_P\in \mathbb{N}$ where $R_F$ is even.
		We call $R_F$ the \emph{external} rounds and $R_P$ the  \emph{internal} rounds: external rounds apply the S-box
		to every lane before mixing, whereas internal rounds apply it only to the first lane. Furthermore, let
		 $\{c^{(i)}_j\}_{i\in [R_F]}^{j\in [t]}\subset\Fq$, $\{\hat c^{(i)}\}_{i\in [R_P]}\subset\Fq$ be public round constants and $\mM_E,\mM_I\in\Fq^{t\times t}$ be invertible maps.
	
	For $\vecx=(x_1,\dots,x_t)\in\Fq^t$, define external rounds
	\[
	E_i(\vecx) := \mM_E \cdot \big((x_1+c^{(i)}_1)^d,\dots,(x_t+c^{(i)}_t)^d\big)^\top,
	\]
	and internal rounds
	\[
	I_i(\vecx) := \mM_I \cdot \big((x_1+\hat c^{(i)})^d,\ x_2,\dots,x_t\big)^\top.
	\]
	Then
	\[
	P := E_{R_F}\circ\cdots\circ E_{R_F/2+1}\circ I_{R_P}\circ\cdots\circ I_{1}\circ E_{R_F/2}\circ\cdots\circ E_1.
	\]
\end{definition}

We additionally define
$\Tr_i:\Fq^t\to \Fq^{i}$ outputs the first $i$ elements of the input vector.

\begin{definition}[Poseidon2 PRF]\label{def:poseidon-prf}
Let $\ell_{\mathsf{tag}},\ell_{\mathsf{key}},\ell_{\mathsf{in}}\in \mathbb{N}$ and set $t=\ell_{\mathsf{key}}+\ell_{\mathsf{in}}$. For a key 
 $\veck\in\Fq^{\ell_{\mathsf{key}}}$  and public input 
$\vecn \in\Fq^{\ell_{\mathsf{in}}}$, we define the Poseidon2 PRF as 
\[
  \PRF(\veck,\vecn) := \Tr_{\ell_{\mathsf{tag}}}\big(P(\veck\Vert \vecn) + (\veck\Vert \vecn)\big),
\]
where $P:\Fq^t\to \Fq^t$ is the permutation from Definition \ref{def:poseidon-perm}.
\end{definition}
